\clearpage
\begingroup
\fancyhead[L]{\small THEORIE DER BEWEGUNG DES MONDES.}
\renewcommand{\headrulewidth}{0pt}
\noindent Ferner findet sich leicht aus der oben gefundenen Formel für \(\theta\):
\noindent\hfill80\textsuperscript{*}\par
\SourceLeaf{642}
\[
\begin{aligned}
(1+\theta\theta)^{-\frac32}
={}&1-\tfrac34gg-\tfrac{27}{256}nngg+\tfrac{45}{64}g^4\\
&+(\tfrac34gg-\tfrac{15}{16}g^4)\cos2\lambda
 +\tfrac{15}{64}g^4\cos4\lambda\\
&-\tfrac{9}{16}ngg(1+\tfrac14n)\cos(2E-2\lambda)\\
&-\tfrac{15}{16}ggee\cos(2M-2\lambda)
 -\tfrac{27}{16}ngg\varepsilon\cos(2\lambda+a)
 +\tfrac{27}{16}ngg\varepsilon\cos(2\lambda-a)\\
&-\tfrac{9}{16}ngg\varepsilon\cos(2E-2\lambda+a)
 +\tfrac{21}{16}ngg\varepsilon\cos(2E-2\lambda-a)\\
&+\tfrac{9}{16}ngg(1+\tfrac14n)\cos2E
 +\tfrac{15}{16}ggee\cos2M\\
&+\tfrac{9}{16}ngg\varepsilon\cos(2E+a)
 -\tfrac{21}{16}ngg\varepsilon\cos(2E-a)
 \bigl[-\tfrac{27}{256}nngg\cos(4E-2\lambda\bigr].
\end{aligned}
\]
\noindent Durch Multiplication dieser Theile und Hinzusetzung des Factors \(k\) findet
man \(p^4\dfrac{\mathrm du^2}{\mathrm dv^2}\dfrac{k}{(1+\theta\theta)^{\frac32}}\).
\par\noindent\hspace*{\parindent}Es bleibt also noch zu entwickeln
\[
\begin{aligned}
&nn\frac{B^3}{R^3}\frac{\mathrm du^2}{\mathrm dv^2}
 \Bigl\{\Bigl(\tfrac12+\tfrac32\cos2(v-V)
 +7\bigl(\tfrac98\cos E+\tfrac{15}{8}\cos3E\bigr)\Bigr)p\\
&\qquad-\Bigl(\tfrac32\sin2(v-V)
 +7\bigl(\tfrac38\cos E+\tfrac{15}{8}\cos3E\bigr)\Bigr)
 \frac{\mathrm dp}{\mathrm dv}\Bigr\}\\
&=nnh\frac{B^3}{R^3}\frac{\mathrm du^2}{\mathrm dv^2}
 \Bigl\{\tfrac12\frac{p}{h}+\tfrac98\,7\cos E+\tfrac{15}{8}\,7\cos3E
 +\tfrac32\cos2(v-V)\\
&\qquad-\tfrac32e\cos(2(v-V)-M)-\tfrac{45}{16}ne\cos M
 +\tfrac94nn-\tfrac34nn\cos4E\\
&\qquad-\tfrac{9}{16}gg\cos(2E-2\lambda)
 +\tfrac{3}{16}gg\cos(2E+2\lambda)\Bigr\}.
\end{aligned}
\]
\noindent Die Entwickelung des letzten Factors gibt
\[
\begin{aligned}
\cos2(v-V)
={}&(1-4\varepsilon\varepsilon)\cos2E
 +2\varepsilon\cos(2E+a)-2\varepsilon\cos(2E-a)\\
&+\tfrac34\varepsilon\varepsilon\cos(2E+2a)
 +\tfrac{13}{4}\varepsilon\varepsilon\cos(2E-2a)\\
&-2ne\cos(2E+M)+2ne\cos(2E-M)\\
\cos(2(v-V)\mathord{.}-M)
={}&\cos(2E-M)+2\varepsilon\cos(2E-M+a)
 -2\varepsilon\cos(2E-M-a).
\end{aligned}
\]
\noindent Folglich
\[
\begin{aligned}
&\frac{\mathrm du^2}{\mathrm dv^2}
 \Bigl(\psi p-\pi\frac{\mathrm dp}{\mathrm dv}\Bigr)=\\
&nnh\frac{B^3}{R^3}\frac{\mathrm du^2}{\mathrm dv^2}
 \Bigl\{\tfrac12\frac{p}{h}+\tfrac94nn
 +\tfrac32(1-4\varepsilon\varepsilon)\cos2E-\tfrac34nn\cos4E
 +\tfrac98\,7\cos E+\tfrac{15}{8}\,7\cos3E\\
&\qquad-\tfrac{45}{16}ne\cos M-3ne\cos(2E+M)
 -\tfrac32e(1-2n)\cos(2E-M)\\
&\qquad+3\varepsilon\cos(2E+a)-3\varepsilon\cos(2E-a)
 +\tfrac98\varepsilon\varepsilon\cos(2E+2a)
 +\tfrac{39}{8}\varepsilon\varepsilon\cos(2E-2a)\\
&\qquad-3e\varepsilon\cos(2E-M+a)+3e\varepsilon\cos(2E-M-a)
 +\tfrac{3}{16}gg\cos(2E+2\lambda)\\
&\qquad-\tfrac{9}{16}gg\cos(2E-2\lambda\Bigr\}.
\end{aligned}
\]
\clearpage
\endgroup
